\documentclass[12pt, a4paper]{article}

\usepackage[utf8]{inputenc}
\usepackage[T1]{fontenc}
\usepackage[spanish]{babel}
\usepackage{geometry}
\usepackage{setspace}
\usepackage{titlesec}
\usepackage{parskip}
\usepackage{fancyhdr}
\usepackage{hyperref}

\geometry{
  top=2.5cm,
  bottom=2.5cm,
  left=3cm,
  right=2.5cm
}

\onehalfspacing

\pagestyle{fancy}
\fancyhf{}
\rhead{\small Extensiones del Lenguaje HULK}
\lhead{\small Universidad de La Habana}
\cfoot{\thepage}
\renewcommand{\headrulewidth}{0.4pt}

\titleformat{\section}{\large\bfseries}{{\thesection.}}{0.5em}{}[\vspace{0.1em}\hrule\vspace{0.4em}]
\titleformat{\subsection}{\normalsize\bfseries}{{\thesubsection.}}{0.5em}{}
\titleformat{\subsubsection}{\normalsize\itshape}{{\thesubsubsection.}}{0.5em}{}

\hypersetup{
  colorlinks=true,
  linkcolor=black,
  urlcolor=black,
  pdfauthor={Compilador HULK},
  pdftitle={Extensiones del Lenguaje HULK},
}

\begin{document}

% ============================================================
% PORTADA
% ============================================================
\begin{titlepage}
  \centering
  \vspace*{2cm}
  {\LARGE\bfseries Informe técnico de lenguaje de programación\par}
  \vspace{0.5cm}
  {\large Extensiones de control de flujo e iteración para HULK\par}
  \vspace{2cm}

  {\large Universidad de La Habana\par}
  \vspace{1cm}
  {\normalsize Facultad de Matemática y Computación\par}
  \vspace{0.5cm}
  {\normalsize \textbf{Integrantes:}\par}
  {\normalsize Daniela de la Caridad Guerrero Alvarez\par}
  {\normalsize Ruben Martinez Rojas\par}
  {\normalsize Sammy Raul Sosa Justiz\par}
  \vspace{0.3cm}

  \vfill
  {\large \today\par}
\end{titlepage}

% ============================================================
% ÍNDICE
% ============================================================
\tableofcontents
\newpage

% ============================================================
% 1. INTRODUCCIÓN
% ============================================================
\section{Introducción y objetivos}

El lenguaje HULK (Havana University Language for Kompilers) fue concebido como un lenguaje de programación didáctico, orientado a objetos, de tipado estático con inferencia de tipos opcional e incremental por diseño. Su naturaleza incremental significa que cada nueva característica puede añadirse sobre una base ya funcional sin romper el subconjunto previamente implementado. Esta filosofía es precisamente la que hace de HULK un campo fértil para la exploración de extensiones sintácticas y semánticas.

El presente informe documenta cuatro extensiones incorporadas al compilador de HULK con el objetivo de ampliar su expresividad en el área del control de flujo y la iteración. Dichas extensiones son: el bucle \texttt{for} tipado con soporte explícito de iteradores, la construcción \texttt{repeat~(n)~expr} para repetición acotada, la negación condicional mediante \texttt{unless~(cond)~expr}, y el bucle postcondicional \texttt{loop~expr~while~(cond)}. Estas cuatro construcciones no son ideas aisladas: conforman un paquete coherente de mejoras orientadas a reducir la fricción entre la intención del programador y el código que debe escribir.

Los objetivos de este informe son tres. Primero, describir cada extensión en detalle: su motivación, su sintaxis propuesta, los cambios que introdujo en la gramática y en el árbol sintáctico abstracto, y su semántica formal mediante desazucarado. Segundo, justificar las decisiones de diseño desde una perspectiva tanto técnica como didáctica. Tercero, comparar cada extensión con construcciones análogas en lenguajes de programación ampliamente conocidos, para contextualizar las elecciones realizadas dentro del panorama más amplio del diseño de lenguajes.

% ============================================================
% 2. ESTADO INICIAL DEL COMPILADOR
% ============================================================
\section{Estado inicial del compilador}

Antes de introducir las extensiones, el compilador de HULK contaba con un conjunto base de construcciones de control de flujo. Ya existía la expresión \texttt{if-elif-else} para condicionales, el bucle \texttt{while} para iteración condicionada en cabeza, y una forma de \texttt{for} que iteraba sobre objetos que implementasen el protocolo \texttt{Iterable}, definido por los métodos \texttt{next()} y \texttt{current()}. También estaban implementadas las expresiones \texttt{let-in} para variables léxicamente acotadas, los bloques de expresiones entre llaves, y el sistema de tipos con herencia simple y protocolos.

La gramática existente seguía la forma LL(1), lo cual impone restricciones precisas sobre cómo se pueden añadir nuevas producciones sin introducir ambigüedades o conflictos. Los archivos centrales sobre los que recaen los cambios de cualquier extensión son, fundamentalmente, la gramática formal en su representación textual, el módulo de construcción del árbol sintáctico concreto, el módulo de conversión de ese árbol a un árbol sintáctico abstracto, los nodos de expresión del AST, y el analizador semántico. El hecho de que parte de la infraestructura para el \texttt{for} ya existiese en el compilador fue un factor determinante para seleccionar las extensiones y jerarquizarlas por nivel de esfuerzo.

% ============================================================
% 3. CRITERIOS DE SELECCIÓN
% ============================================================
\section{Criterios para escoger las extensiones}

La selección de las cuatro extensiones no fue arbitraria. Respondió a un conjunto de criterios técnicos y pedagógicos que se describen a continuación.

El primer criterio fue la reutilización de infraestructura existente. Todas las extensiones se traducen, en última instancia, a construcciones que el compilador ya sabe manejar: \texttt{while}, \texttt{if}, \texttt{let}. Esto minimiza el riesgo de introducir errores profundos y permite que cada extensión quede bien acotada en términos de los archivos que toca.

El segundo criterio fue la coherencia temática. Las cuatro extensiones pertenecen al dominio del control de flujo y la iteración. Presentarlas como un paquete unificado, en lugar de cuatro adiciones dispares, refuerza la narrativa de diseño y facilita la escritura del informe académico.

El tercer criterio fue la gradación de complejidad. Se eligió una extensión principal de complejidad media-alta (\texttt{for} tipado con iteradores) y tres extensiones secundarias de complejidad baja que funcionan principalmente como azúcar sintáctica. Esta gradación permite cumplir el requisito mínimo de al menos una extensión con cambios simultáneos en sintaxis y semántica, mientras se aprovecha el trabajo para generar documentación rica.

El cuarto criterio fue la comparabilidad con otros lenguajes. Las cuatro extensiones tienen equivalentes directos en lenguajes bien conocidos, lo que facilita la justificación del diseño y enriquece el análisis comparativo.

% ============================================================
% 4. EXTENSIÓN PRINCIPAL: for TIPADO CON ITERADORES
% ============================================================
\section{Extensión principal: \texttt{for} tipado con iteradores}

\subsection{Motivación}

El lenguaje HULK ya disponía de una forma básica de \texttt{for} que permitía iterar sobre cualquier expresión que implementase el protocolo \texttt{Iterable}. Sin embargo, en esa forma original la variable de iteración no podía recibir una anotación de tipo explícita, lo que en ciertos contextos obligaba al programador a depender completamente de la inferencia de tipos o a recurrir a conversiones manuales mediante el operador \texttt{as}. La extensión propuesta añade la posibilidad de anotar el tipo de la variable de iteración directamente en la cabecera del \texttt{for}, haciendo el código más legible y permitiendo al verificador de tipos detectar inconsistencias de manera más temprana.

La motivación también tiene una dimensión semántica. Al permitir la anotación explícita del tipo de la variable de iteración, el compilador puede verificar estáticamente que el tipo que devuelve \texttt{current()} del iterador sea compatible con el tipo declarado, en lugar de diferir esa verificación o realizarla implícitamente.

\subsection{Sintaxis propuesta}

La forma original del \texttt{for} en HULK tenía la siguiente sintaxis:

\begin{center}
\texttt{for (x in expr) body}
\end{center}

La extensión añade una variante opcional que permite anotar el tipo de la variable de iteración:

\begin{center}
\texttt{for (x : T in expr) body}
\end{center}

Ambas formas conviven en la gramática: la presencia de los dos puntos y el identificador de tipo es opcional. Esto garantiza compatibilidad regresiva total con el código HULK ya existente.

\subsection{Cambios en la gramática}

Para admitir la nueva variante, la producción del \texttt{for} en la gramática LL(1) fue modificada para contemplar un bloque de anotación de tipo optativo inmediatamente después del identificador de la variable de iteración. Dado que la gramática es LL(1), fue necesario asegurarse de que el símbolo de anticipación permitiese distinguir sin ambigüedad la presencia o ausencia de la anotación, lo cual se logró aprovechando que los dos puntos son un terminal que no puede aparecer en ninguna otra posición válida dentro de la cabecera del \texttt{for}.

\subsection{Cambios en el AST}

El nodo del AST que representa el \texttt{for} fue extendido para almacenar un campo adicional optativo correspondiente al tipo anotado de la variable de iteración. Cuando ese campo está ausente, el comportamiento del nodo es exactamente el mismo que antes de la extensión. Cuando está presente, el analizador semántico lo usa para generar una anotación de tipo explícita sobre la variable de iteración durante la fase de desazucarado.

\subsection{Semántica por desazucarado}

La semántica del \texttt{for} tipado se define, como en HULK original, por transpilación a su equivalente en \texttt{while}. La forma sin anotación se desazucara como siempre:

\begin{center}
\texttt{let \_iter = expr in while (\_iter.next()) let x = \_iter.current() in body}
\end{center}

La forma con anotación de tipo \texttt{T} se desazucara de manera similar, pero añadiendo la anotación al \texttt{let} interno:

\begin{center}
\texttt{let \_iter = expr in while (\_iter.next()) let x : T = \_iter.current() in body}
\end{center}

Esto hace que el verificador de tipos compruebe automáticamente que el valor retornado por \texttt{current()} sea compatible con \texttt{T}, sin necesidad de añadir ninguna lógica especial al verificador.

\subsection{Justificación del diseño}

La decisión de mantener la anotación como optativa, en lugar de obligatoria, responde a la filosofía de HULK de permitir un tipado gradual. Un programador que confía en la inferencia de tipos no necesita cambiar nada. Uno que quiere ser explícito o que trabaja con jerarquías de tipos complejas puede hacerlo sin coste adicional. La implementación por desazucarado garantiza además que la extensión no añade complejidad algorítmica al compilador: toda la lógica de verificación de tipos que ya existía para \texttt{let-in} se reutiliza íntegramente.

% ============================================================
% 5. EXTENSIÓN SECUNDARIA: repeat
% ============================================================
\section{Extensión secundaria: \texttt{repeat~(n)~expr}}

\subsection{Motivación}

Uno de los patrones más frecuentes en programación es la repetición de una acción un número fijo y conocido de veces. En HULK base, esto se expresa mediante un bucle \texttt{while} con un contador explícito o mediante un \texttt{for} sobre un rango. Ambas formas son correctas, pero ninguna de ellas expresa directamente la intención del programador cuando simplemente quiere que algo suceda diez veces. La construcción \texttt{repeat~(n)~expr} cubre exactamente ese caso de uso con una sintaxis clara, directa e inmediatamente legible.

\subsection{Sintaxis propuesta}

\begin{center}
\texttt{repeat (n) expr}
\end{center}

Donde \texttt{n} es una expresión de tipo \texttt{Number} que determina el número de repeticiones, y \texttt{expr} es cualquier expresión válida de HULK, incluyendo bloques de expresiones entre llaves. El resultado de la construcción es el último valor evaluado por \texttt{expr}, o \texttt{Null} si \texttt{n} es cero o negativo.

\subsection{Cambios en la gramática y el AST}

Se añadió una nueva producción a la gramática que reconoce la palabra clave \texttt{repeat} seguida de una expresión entre paréntesis y una expresión cuerpo. El nodo correspondiente en el AST almacena la expresión de conteo y la expresión cuerpo. Dado que se trata de una azúcar sintáctica pura, el nodo puede ser eliminado por completo durante la fase de desazucarado, sin dejar rastro en el árbol posterior.

\subsection{Semántica por desazucarado}

La semántica de \texttt{repeat~(n)~expr} se define por su equivalente en \texttt{while}:

\begin{center}
\texttt{let \_total = n in while (\_total > 0) \{ \_total := \_total - 1; expr \}}
\end{center}

El prefijo de guion bajo en \texttt{\_total} sigue la convención de HULK para variables introducidas por el compilador, garantizando que no colisionen con variables del programador. La verificación semántica añade únicamente la comprobación de que \texttt{n} sea de tipo \texttt{Number}, lo que es trivial dado el sistema de tipos existente.

\subsection{Justificación del diseño}

La decisión de desazucarar \texttt{repeat} a \texttt{while} en lugar de a \texttt{for} sobre un rango obedece a que la intención semántica es diferente: \texttt{repeat} no itera sobre elementos, sino que cuenta ejecuciones. Usar \texttt{for} introduciría una variable de iteración que el programador no necesita y que podría provocar confusión. La forma en \texttt{while} con contador descendente es más fiel a la intención de la construcción y genera código intermedio más eficiente en términos conceptuales.

% ============================================================
% 6. EXTENSIÓN SECUNDARIA: unless
% ============================================================
\section{Extensión secundaria: \texttt{unless~(cond)~expr}}

\subsection{Motivación}

La construcción condicional negada es un patrón que aparece con frecuencia en código imperativo: ejecutar algo si y solo si una condición \emph{no} se cumple. La forma canónica en HULK sería \texttt{if~(!cond)~expr}, pero cuando la condición es larga o compleja, la negación explícita con \texttt{!} reduce la legibilidad y obliga al programador a hacer un salto mental para entender la intención del código. La construcción \texttt{unless} elimina ese salto: el código dice exactamente lo que hace, sin negaciones que el lector deba mentalmente invertir.

\subsection{Sintaxis propuesta}

La forma base de la extensión es la siguiente:

\begin{center}
\texttt{unless (cond) expr}
\end{center}

Y admite también una variante con rama alternativa:

\begin{center}
\texttt{unless (cond) expr else alt}
\end{center}

La variante con \texttt{else} permite que la construcción sea usada como expresión con valor en ambas ramas, tal como ocurre con \texttt{if} en HULK.

\subsection{Cambios en la gramática y el AST}

La palabra clave \texttt{unless} fue añadida al léxico del lenguaje y reconocida como un nuevo terminal. En la gramática, se añadió una producción que sigue la misma estructura que la de \texttt{if}, con la diferencia de que no admite ramas \texttt{elif}: solo la forma directa y la variante con \texttt{else}. Esta restricción es intencional, ya que \texttt{unless} está diseñado para casos de uso simples; cuando la lógica se vuelve compleja, el programador debería usar \texttt{if} con sus múltiples ramas. El nodo del AST para \texttt{unless} almacena la condición y hasta dos ramas opcionales.

\subsection{Semántica por desazucarado}

El desazucarado de \texttt{unless} es directo. La forma sin \texttt{else}:

\begin{center}
\texttt{unless (cond) expr}  $\;\Rightarrow\;$  \texttt{if (!cond) expr}
\end{center}

Y la forma con \texttt{else}:

\begin{center}
\texttt{unless (cond) expr else alt}  $\;\Rightarrow\;$  \texttt{if (!cond) expr else alt}
\end{center}

Tras el desazucarado, toda la verificación semántica es heredada de \texttt{if}, incluyendo la comprobación de que la condición sea de tipo booleano y la inferencia del tipo de resultado como el ancestro común más bajo de los tipos de ambas ramas.

\subsection{Justificación del diseño}

La decisión más relevante en el diseño de \texttt{unless} fue la de no admitir \texttt{elif}. Esta restricción evita que \texttt{unless} se convierta en un duplicado complejo de \texttt{if} con lógica invertida, lo cual sería confuso. La construcción se mantiene simple y de propósito único: invertir una condición simple. Cualquier uso más elaborado de lógica condicional sigue siendo territorio de \texttt{if}.

% ============================================================
% 7. EXTENSIÓN SECUNDARIA: loop ... while
% ============================================================
\section{Extensión secundaria: \texttt{loop~expr~while~(cond)}}

\subsection{Motivación}

Los bucles \texttt{while} tradicionales evalúan la condición antes de ejecutar el cuerpo, lo que significa que si la condición es falsa desde el inicio, el cuerpo nunca se ejecuta. Hay, sin embargo, situaciones en las que se desea ejecutar el cuerpo al menos una vez antes de comprobar la condición. Este patrón, conocido como bucle postcondicional, es común en la lectura de entradas del usuario, en el procesamiento de flujos de datos, y en algoritmos que siempre realizan al menos una iteración. La construcción \texttt{loop~expr~while~(cond)} cubre ese caso de uso de forma clara y declarativa.

\subsection{Sintaxis propuesta}

\begin{center}
\texttt{loop expr while (cond)}
\end{center}

Donde \texttt{expr} es la expresión cuerpo que se ejecuta al menos una vez, y \texttt{cond} es la condición que determina si se continúa iterando. El resultado de la construcción es el último valor evaluado por \texttt{expr}.

\subsection{Cambios en la gramática y el AST}

Se añadió la palabra clave \texttt{loop} como nuevo terminal. La producción correspondiente en la gramática reconoce \texttt{loop}, seguido de la expresión cuerpo, la palabra clave \texttt{while} y la condición entre paréntesis. El nodo del AST para esta construcción almacena la expresión cuerpo y la condición, en ese orden, reflejando la semántica de ejecución: primero el cuerpo, luego la evaluación de la condición.

\subsection{Semántica por desazucarado}

La semántica de \texttt{loop~expr~while~(cond)} se define por el siguiente desazucarado:

\begin{center}
\texttt{let \_done = false in while (!\_done) \{ expr; \_done := !(cond); \}}
\end{center}

Este desazucarado garantiza que \texttt{expr} se ejecuta exactamente una vez antes de que \texttt{cond} sea evaluado por primera vez. Si \texttt{cond} resulta verdadero tras la primera ejecución, el bucle continúa. Si resulta falso, el bucle termina. Al igual que con \texttt{repeat}, la variable auxiliar lleva el prefijo de guion bajo para evitar colisiones.

\subsection{Justificación del diseño}

La elección de la palabra clave \texttt{loop} en lugar de la forma \texttt{do~...~while} de C y Java fue deliberada. La forma \texttt{do~...~while} introduce dos palabras clave (\texttt{do} y el \texttt{while} final) y rompe la convención de que en HULK los bloques se identifican visualmente por las llaves o por la posición del cuerpo respecto a la condición. La forma \texttt{loop~...~while} es más uniforme con la gramática existente y sitúa la condición al final de forma que el lector sabe desde el inicio que está ante un bucle cuyo cuerpo precede a su condición de corte.

% ============================================================
% 8. COMPARATIVA CON OTROS LENGUAJES
% ============================================================
\section{Comparativa con otros lenguajes de programación}

\subsection{El bucle \texttt{for} con anotación de tipo}

La extensión del \texttt{for} en HULK encuentra sus análogos más directos en varios lenguajes modernos. En Python, el bucle \texttt{for~x~in~items} no lleva anotación de tipo en la variable de iteración por defecto, aunque el sistema de anotaciones de tipo de Python 3 permite declarar el tipo de una variable mediante un comentario o una anotación separada; sin embargo, esto nunca forma parte de la sintaxis del \texttt{for} mismo. En C\#, el equivalente es \texttt{foreach~(T~x~in~collection)}, donde el tipo de la variable de iteración es parte integral de la sintaxis, lo que lo hace el precedente más cercano a la extensión implementada en HULK. En Kotlin, la forma es \texttt{for~(x:T~in~xs)}, con la anotación de tipo separada por dos puntos exactamente como en la extensión de HULK, lo que no es una coincidencia sino una decisión consciente de adoptar una sintaxis que ya ha demostrado ser legible y popular. En Rust, el bucle \texttt{for~x~in~iter} tampoco admite anotaciones de tipo directamente en la cabecera, pero el compilador infiere el tipo con mucha precisión gracias a su sistema de tipos basado en rasgos, que es conceptualmente similar al sistema de protocolos de HULK.

La principal diferencia entre la extensión de HULK y sus equivalentes es la optatividad de la anotación. Mientras que en C\# la anotación de tipo en \texttt{foreach} es obligatoria (aunque puede usarse la palabra clave \texttt{var} para inferencia), en HULK la anotación sigue siendo totalmente opcional, lo cual preserva la flexibilidad del tipado gradual.

\subsection{La construcción \texttt{repeat}}

La construcción \texttt{repeat~(n)~expr} tiene precedentes en varios lenguajes orientados a la expresividad y la concisión. En Ruby, la forma \texttt{n.times~\{~bloque~\}} es idiomática para repetir un bloque un número fijo de veces, y se usa ampliamente en código Ruby para tareas simples de repetición. En Lua, el bucle \texttt{repeat~...~until} existe, aunque con semántica diferente: en Lua es postcondicional y termina cuando la condición es verdadera, que es lo opuesto del \texttt{while}. En Pascal, el bucle \texttt{for} estándar cubre casos de repetición acotada con más verbosidad. En lenguajes de dominio específico orientados a la educación como Logo o Scratch, la construcción de repetición directa es prácticamente ubicua y suele ser la primera estructura de control que los estudiantes aprenden.

La diferencia más notable entre la construcción de HULK y el \texttt{n.times} de Ruby es de naturaleza sintáctica: en Ruby, \texttt{times} es un método del objeto numérico, lo que implica que los enteros son objetos con métodos, una decisión de diseño que HULK no ha adoptado. En HULK, \texttt{repeat} es una construcción de lenguaje de primer orden, más parecida a una palabra clave que a una llamada a método, lo que la hace más natural dentro del paradigma del lenguaje.

\subsection{La construcción \texttt{unless}}

La negación condicional tiene su precedente más famoso en Ruby, donde \texttt{unless~cond~do~...~end} es una construcción idiomática ampliamente usada. Ruby fue probablemente el primer lenguaje mainstream en adoptar \texttt{unless} como ciudadano de primera clase, y su popularidad entre los programadores Ruby es un testimonio de que la construcción mejora genuinamente la legibilidad en los casos de uso para los que fue diseñada.

En Swift, los \emph{guard statements} cumplen una función similar pero más poderosa: permiten salir anticipadamente de una función si una condición no se cumple. En Kotlin, las expresiones de guarda tienen un rol análogo. Estas construcciones son más expresivas que un simple \texttt{unless}, pero también más complejas. La extensión de HULK es deliberadamente más modesta: solo invierte la condición y no introduce semántica de salida anticipada, lo cual es coherente con la naturaleza de HULK como lenguaje basado en expresiones con un único valor de retorno por bloque.

En Perl, \texttt{unless} también existe y puede usarse tanto en forma prefija como en forma postfija. La versión postfija, como \texttt{print~(x)~unless~(x~==~0)}, está ausente de la extensión de HULK por razones gramaticales: la forma postfija requeriría que el cuerpo de la expresión aparezca antes que la palabra clave, lo cual introduce dificultades en una gramática LL(1).

\subsection{La construcción \texttt{loop ... while}}

El bucle postcondicional tiene su equivalente más conocido en el \texttt{do~...~while} de C, C++ y Java. Esta construcción es una de las más antiguas del paradigma imperativo y su utilidad está bien establecida. La diferencia fundamental entre \texttt{do~...~while} y la construcción \texttt{loop~...~while} de HULK es, como ya se mencionó, de orden léxico y de coherencia visual con el resto del lenguaje.

En Pascal, la construcción análoga es \texttt{repeat~...~until}, donde la condición de terminación es la negación lógica de la condición de continuación: el bucle termina cuando la condición es verdadera, en lugar de cuando es falsa. Esta inversión puede resultar confusa para programadores acostumbrados al paradigma de C, por lo que HULK optó por conservar la semántica de \texttt{while}: el bucle continúa mientras la condición sea verdadera, lo cual es consistente con el bucle \texttt{while} precondicional que ya existía en el lenguaje.

En algunos lenguajes funcionales con características imperativas, como Scala o Groovy, los bucles postcondicionales son menos comunes porque el paradigma funcional favorece la recursión y las comprensiones de listas sobre la iteración explícita. HULK, siendo un lenguaje orientado a objetos con raíces imperativas claras, encaja más naturalmente con la construcción postcondicional de la tradición C.

% ============================================================
% 9. DECISIONES DE DISEÑO Y TRADEOFFS
% ============================================================
\section{Decisiones de diseño y tradeoffs}

\subsection{Azúcar sintáctica versus extensión semántica}

La tensión más importante en el diseño de estas extensiones es la que existe entre la riqueza semántica y la simplicidad de implementación. Las tres extensiones secundarias son azúcar sintáctica pura: no añaden ninguna capacidad computacional que no existiese antes. El \texttt{for} tipado sí añade una capacidad verificadora nueva, pero incluso esta se implementa por desazucarado a construcciones existentes.

Esta decisión tiene consecuencias positivas y negativas. En el lado positivo, el riesgo de introducir errores es mínimo, la carga de mantenimiento es baja, y toda la infraestructura de generación de código intermedio y de optimización que ya existía sigue funcionando sin modificaciones. En el lado negativo, las extensiones no pueden hacer nada que no pudiese hacerse antes, lo que implica que su valor es puramente expresivo y no computacional.

La justificación de este tradeoff es precisamente que en un lenguaje didáctico, la expresividad tiene un valor enorme. Un estudiante que aprende a programar con HULK debería poder escribir código que se lea como prosa, y las extensiones contribuyen a ese objetivo sin añadir complejidad al modelo computacional subyacente.

\subsection{Compatibilidad regresiva}

Todas las extensiones se diseñaron para ser completamente compatibles con el código HULK existente. Ninguna palabra clave nueva colisiona con identificadores válidos en HULK base, y ninguna extensión modifica la semántica de las construcciones existentes. Esto garantiza que cualquier programa escrito para el compilador base sigue compilando y ejecutándose correctamente después de añadir las extensiones.

\subsection{Uniformidad sintáctica}

Un principio rector en el diseño de las extensiones fue la uniformidad sintáctica. Las nuevas construcciones siguen los mismos patrones visuales que las existentes: las condiciones van entre paréntesis, los cuerpos son expresiones que pueden ser bloques entre llaves, y las palabras clave son minúsculas. Esto hace que un programador familiarizado con HULK pueda leer código que usa las extensiones sin necesidad de consultar documentación adicional.

\subsection{El problema de la expresividad y la sorpresa}

Una consideración importante fue evitar lo que en diseño de lenguajes se llama el principio de la menor sorpresa. Cada construcción hace exactamente lo que un programador razonable esperaría que hiciese. \texttt{unless} invierte la condición. \texttt{repeat} repite. \texttt{loop~...~while} ejecuta el cuerpo antes de evaluar la condición. Ninguna de estas construcciones tiene efectos secundarios implícitos ni comportamientos ocultos que puedan sorprender al programador.

% ============================================================
% 10. CAMBIOS EN GRAMÁTICA, AST Y SEMÁNTICA
% ============================================================
\section{Resumen de cambios en gramática, AST y semántica}

Los cambios en la gramática se concentraron en añadir nuevas producciones para cada una de las cuatro construcciones. Para el \texttt{for} tipado, se modificó la producción existente para hacer opcional la anotación de tipo. Para \texttt{repeat}, \texttt{unless} y \texttt{loop~...~while} se añadieron producciones completamente nuevas. En todos los casos se verificó que la gramática resultante siguiese siendo LL(1), comprobando la ausencia de ambigüedades y conflictos en los conjuntos de predicción.

Los cambios en el AST consistieron en añadir un nuevo nodo por cada construcción nueva, y en extender el nodo existente del \texttt{for} con un campo optativo para la anotación de tipo. Cada nuevo nodo almacena exactamente la información necesaria para el desazucarado: la condición (cuando la hay), la expresión de conteo (para \texttt{repeat}), y el cuerpo.

Los cambios en el analizador semántico se limitaron a añadir las verificaciones específicas de cada construcción antes del desazucarado: que la condición de \texttt{unless} sea booleana, que el argumento de \texttt{repeat} sea numérico, y que la condición de \texttt{loop~...~while} sea booleana. Una vez superadas esas verificaciones, el desazucarado convierte cada nueva construcción en su equivalente base, y la verificación semántica completa se delega al sistema ya existente.

% ============================================================
% 11. CASOS DE PRUEBA
% ============================================================
\section{Casos de prueba}

La validación de las extensiones se realizó mediante un conjunto de casos de prueba que cubren tanto los casos de uso esperados como los casos límite.

Para el \texttt{for} tipado, se probó la iteración sobre rangos numéricos con y sin anotación de tipo, la iteración sobre colecciones personalizadas que implementan el protocolo \texttt{Iterable}, y los casos de error en los que el tipo anotado no es compatible con el tipo devuelto por \texttt{current()}.

Para \texttt{repeat}, se probó la repetición de un número positivo de veces, el comportamiento con un argumento cero y con un argumento negativo, y la anidación de \texttt{repeat} dentro de otras construcciones de control de flujo.

Para \texttt{unless}, se probó la forma sin \texttt{else} con una condición verdadera (donde el cuerpo no debe ejecutarse) y con una condición falsa (donde sí debe ejecutarse), así como la forma con \texttt{else} en ambas situaciones, verificando que el tipo de retorno se infiere correctamente como el ancestro común más bajo de los tipos de ambas ramas.

Para \texttt{loop~...~while}, se probó el caso en que la condición es falsa desde el inicio pero el cuerpo se ejecuta una vez, el caso de múltiples iteraciones, y la anidación con otras construcciones.

% ============================================================
% 12. LIMITACIONES Y TRABAJO FUTURO
% ============================================================
\section{Limitaciones y trabajo futuro}

Las extensiones presentadas tienen varias limitaciones que podrían abordarse en iteraciones futuras del compilador.

La construcción \texttt{repeat} no expone la variable de contador al programador. En lenguajes como Ruby, \texttt{n.times} permite recibir el índice actual en el bloque si el programador lo desea. Añadir esa capacidad a \texttt{repeat} en HULK requeriría introducir un placeholder de variable en el estilo de los macros del lenguaje, lo que aumentaría la complejidad de la implementación pero mejoraría significativamente la utilidad de la construcción.

La construcción \texttt{unless} no admite ramas \texttt{elif}. Si bien esta restricción fue deliberada, podría ser reconsiderada si el análisis de uso mostrara que los programadores frecuentemente quieren encadenar condiciones negadas.

El bucle \texttt{loop~...~while} no tiene una forma de salida anticipada equivalente a \texttt{break}. HULK en general no tiene \texttt{break} ni \texttt{continue}, lo que es una limitación del lenguaje base más que de la extensión. Añadir mecanismos de salida no local es una extensión de mayor envergadura que podría considerarse en el futuro.

El \texttt{for} tipado no verifica en tiempo de compilación que el tipo anotado sea exactamente el tipo de \texttt{current()}, sino que verifica conformidad. Esto es correcto según la semántica de HULK, pero podría ser sorprendente para programadores acostumbrados a lenguajes con tipado más rígido.

% ============================================================
% 13. CONCLUSIONES
% ============================================================
\section{Conclusiones}

Las cuatro extensiones implementadas, el \texttt{for} tipado con iteradores, la construcción \texttt{repeat}, la negación condicional \texttt{unless} y el bucle postcondicional \texttt{loop~...~while}, conforman un paquete coherente de mejoras al lenguaje HULK en el dominio del control de flujo y la iteración.

Desde el punto de vista técnico, las extensiones demuestran cómo es posible enriquecer significativamente la expresividad de un lenguaje sin añadir complejidad al modelo computacional subyacente. El mecanismo del desazucarado, que transpila cada nueva construcción a su equivalente en construcciones base, resulta ser una estrategia de implementación robusta, mantenible y segura. No requiere cambios en la generación de código intermedio, no introduce nuevas fases al compilador, y delega toda la verificación semántica compleja al sistema ya existente.

Desde el punto de vista del diseño de lenguajes, las extensiones ilustran varios principios fundamentales. La compatibilidad regresiva garantiza que las mejoras no rompan el código existente. La uniformidad sintáctica hace que las nuevas construcciones sean inmediatamente legibles para cualquier programador familiarizado con el lenguaje. La gradación de complejidad permite que las extensiones sean implementadas de forma incremental, comenzando por las más simples y avanzando hacia las que requieren más cambios en el sistema de tipos.

La comparativa con otros lenguajes revela que ninguna de las cuatro extensiones es original en el panorama del diseño de lenguajes: todas tienen precedentes bien establecidos en Python, Ruby, Kotlin, C, C\#, Pascal y otros. Lo que sí es original es la forma en que se integran coherentemente en la filosofía de HULK, respetando su modelo de expresiones, su sistema de tipos gradual y su gramática LL(1).

En última instancia, estas extensiones son una demostración de que HULK, como lenguaje diseñado para crecer de forma incremental, está bien equipado para acoger nuevas características sin perder su identidad ni su coherencia interna.

\end{document}
